\section{結論}

本稿では，日本語の母音列のみを入力情報とし，
LoRA により追加学習した小型 LLM を用いて
文脈に応じた日本語候補を生成・順位付けする超省入力型日本語入力システムを提案した．
従来 IME の変換精度ではなく入力操作量に着目し，
入力を 5 母音と撥音にまで削減したうえで，
失われた情報を LLM の文脈推定で補う枠組みを示した．
実装としては，CC-100 と字幕コーパスから母音列$\to$文の SFT データを構築し，
MiniCPM4-0.5B を LoRA で追加学習し，ビーム探索により上位 $k$ 候補を提示する系を構築した．
評価については，候補生成精度（Acc@$k$, MRR, CER/WER），
入力効率（KSPC, 入力削減率），UI 効率（キー数，画面占有率），
および認知負荷（NASA-TLX を参考にしつつ，母音変換負荷と記憶保持負荷を
独立に測る主観項目）を従来 IME と比較する枠組みを整備した．

今後の課題は，(1) 評価コードによる定量指標の確定，
(2) Base LLM・頻度 baseline・random baseline との比較実装，
(3) 若年・高齢，フリック熟練・非熟練を含む実ユーザ実験の実施，
(4) 固有名詞・数字・英字に対する補完経路の設計である．
